% =====================================================================
\lezione{02/11/2026}{Conclusione CRA e simulazione CRA--GDPR (Pulse K)}
  {Prof.~Giuseppe Vaciago; Gianluca (laboratori)}
  {obblighi di segnalazione CRA; CRA e AI Act; impatti del CRA; software nel perimetro CRA;
   simulazione di incidente Pulse K; presentazione del laboratorio di audit}
\label{lez:2026-11-02}
% =====================================================================

\begin{nota}[Fonte]
Trascrizione automatica della registrazione, revisionata. Le date, gli orari e i numeri del caso
Pulse K sono ripresi dalle slide \emph{Case1\_PulseK\_CRA\_GDPR\_Simulation} (la trascrizione era
inaffidabile su questi punti). Slide di riferimento: \emph{05.~PoliTO\_CRA}.
\end{nota}

\section{Sintesi}
Conclusione della parte teorica sul Cyber Resilience Act: obblighi di segnalazione, rapporto con
l'AI Act, impatti e sfide. Segue la prima simulazione di incidente (smartwatch per bambini Pulse K),
che obbliga a lavorare insieme su CRA e GDPR: qualificazione del caso, orologi di notifica,
severità ENISA e piano di remediation. Messaggio centrale: a fine corso si deve saper
\textbf{(1) riconoscere la norma applicabile} e \textbf{(2) seguire il percorso corretto in caso di
incidente}. Infine Gianluca presenta il laboratorio di venerdì, in cui gli studenti faranno gli auditor.

% ---------------------------------------------------------------------
\section{Obblighi di segnalazione del CRA: riepilogo}

Il fabbricante (\emph{manufacturer}) deve notificare ogni \textbf{vulnerabilità attivamente
sfruttata} presente in un prodotto con elementi digitali di cui venga a conoscenza. La notifica
passa dalla \textbf{Single Reporting Platform} gestita da ENISA e raggiunge \textbf{due destinatari}:
\begin{itemize}
  \item il \textbf{CSIRT} designato come coordinatore (per l'Italia il CSIRT Italia presso l'ACN);
  \item l'\textbf{ENISA}, che ha accesso alla piattaforma.
\end{itemize}

\paragraph{Vulnerabilità attivamente sfruttata (art.~14(2)).}
\begin{itemize}
  \item \textbf{Early warning}: entro \textbf{24 ore} dalla conoscenza;
  \item \textbf{Notifica della vulnerabilità}: senza ingiustificato ritardo e comunque entro \textbf{72 ore};
  \item \textbf{Final report}: entro \textbf{14 giorni} dal momento in cui è disponibile una misura
    correttiva o di mitigazione.
\end{itemize}

\paragraph{Incidente grave (art.~14(4)).} Lo schema è parallelo: early warning entro 24 ore,
notifica dell'incidente entro 72 ore, \textbf{final report entro un mese}. Il professore non ha
messo questa scadenza nelle slide per non aggiungere complessità, ma ha anticipato che comparirà
nella tabella riassuntiva preparata per l'esame.

\paragraph{Date di applicazione.} Obblighi di segnalazione dall'\textbf{11 settembre 2026};
applicazione integrale dall'\textbf{11 dicembre 2027}.

\paragraph{Differenza principale con il GDPR.}
Il CRA fissa una scadenza anche per il \emph{final report}, il GDPR no. Nel GDPR entro 72 ore va
fatta la notifica, ma raccogliere tutte le informazioni del modulo in 72 ore è quasi impossibile:
nella prassi la notifica completa arriva dopo almeno un mese e, trascorse due o tre settimane, è
il Garante a sollecitarne il completamento. Non esiste però un termine vincolante.

\begin{esame}
\begin{itemize}
  \item Primo passo: \textbf{qualificare il caso} (CRA, NIS2, DORA, GDPR?).
  \item Secondo passo: ricordare con precisione le \textbf{tempistiche} applicabili
    (vedi tabella al Capitolo~\ref{cap:scadenze}).
  \item Errore tipico: confondere l'\emph{early warning} del CRA con la \emph{notifica preliminare}
    del GDPR.
\end{itemize}
\end{esame}

% ---------------------------------------------------------------------
\section{CRA e AI Act: presunzione di conformità}

I prodotti con elementi digitali che rientrano nel CRA e sono anche \textbf{sistemi di IA ad alto
rischio} ai sensi dell'AI Act si considerano conformi ai requisiti di cybersicurezza dell'AI Act
(art.~15 AI Act) se sono soddisfatte tre condizioni:
\begin{enumerate}
  \item il prodotto soddisfa i requisiti essenziali dell'\textbf{Allegato I, Parte I} del CRA;
  \item i processi del fabbricante rispettano l'\textbf{Allegato I, Parte II} del CRA
    (gestione delle vulnerabilità);
  \item il livello di protezione richiesto dall'AI Act è dimostrato nella \textbf{dichiarazione di
    conformità UE} rilasciata ai sensi del CRA.
\end{enumerate}

\noindent Anticipazione sull'AI Act: esistono \textbf{pratiche vietate} (es.\ \emph{social scoring})
e \textbf{sistemi ad alto rischio}, cioè sistemi che possono causare un danno serio alle persone e
sono soggetti a misure specifiche, anche di cybersicurezza. L'Allegato III dell'AI Act verrà
approfondito più avanti.

\begin{esame}
Da ricordare: esiste un'interazione forte tra l'\textbf{art.~15 dell'AI Act} e il CRA per i sistemi
di IA ad alto rischio. Gli allegati non vanno imparati a memoria.
\end{esame}

% ---------------------------------------------------------------------
\section{Impatti e sfide del CRA}

Il settore \textbf{tech e sviluppo software} sarà probabilmente il più colpito, data la pervasività
degli elementi digitali: potranno servire cambiamenti significativi nelle pratiche di sviluppo e nei
prodotti. Le sfide principali sono:
\begin{itemize}
  \item complessità dei requisiti;
  \item integrazione della \emph{security by design} e \emph{by default} in prodotti e processi;
  \item gestione della \emph{supply chain};
  \item investimenti necessari;
  \item gestione continua del rischio.
\end{itemize}

% ---------------------------------------------------------------------
\section{Perché il CRA è complesso: il quadro d'insieme}

\begin{itemize}
  \item \textbf{NIS2}: riguarda solo alcune categorie (soggetti \emph{essenziali} e \emph{importanti}),
    legate a infrastrutture critiche, con requisiti di sicurezza specifici e numerosi.
  \item \textbf{GDPR}: si applica a chiunque tratti dati personali, senza distinzioni di categoria.
    L'art.~32 discende dal principio di \emph{accountability} e non prescrive regole tecniche puntuali.
  \item \textbf{AI Act}: riguarda in primo luogo i \emph{provider} di sistemi di IA.
  \item \textbf{CRA}: l'esempio classico è lo smartwatch o qualsiasi oggetto IoT, ma il regolamento
    copre anche il \textbf{software}. Un'applicazione può quindi ricadere insieme sotto CRA, NIS2
    (che include il settore digitale) e GDPR.
\end{itemize}

\begin{esame}
\textbf{Le due competenze da portare a casa dal corso:}
\begin{enumerate}
  \item \textbf{ricondurre il rischio alla norma corretta}: capire quale quadro normativo si applica
    a un prodotto, un sistema di IA o una situazione. Non serve diventare giuristi: per i dettagli
    ci si può appoggiare a consulenti esperti della singola norma, ma il perimetro va riconosciuto da soli;
  \item \textbf{in caso di incidente, individuare il percorso corretto} da seguire: a chi notificare,
    cosa e entro quando.
\end{enumerate}
I requisiti tecnici, legali e organizzativi elencati negli allegati non vanno imparati a memoria.
Il problema vero, secondo il professore, è non capire in quale quadro normativo ci si trova.
\end{esame}

% ---------------------------------------------------------------------
\section{Chiarimento: il software nel perimetro del CRA}

Su domanda di uno studente, il professore ha approfondito la nozione di \emph{prodotto con elementi
digitali} applicata al software:
\begin{itemize}
  \item l'elemento chiave è la \textbf{connessione} (logica o fisica) a un dispositivo o a una rete;
  \item un \textbf{software installabile} (un'applicazione) rientra nel CRA; un sito web di per sé no;
  \item rientra anche il software \textbf{connesso a un hardware};
  \item restano esclusi gli ambiti coperti da regimi settoriali (es.\ software sanitario soggetto a
    normativa specifica) e l'ambito militare e della difesa.
\end{itemize}

\begin{esame}
\textbf{Il software rientra nel CRA.} Domande di dettaglio su questa definizione non saranno chieste,
perché la norma non è ancora pienamente applicabile e alcune interpretazioni restano incerte.
\end{esame}

% =====================================================================
\section{Simulazione \#1 -- Pulse K (Volta Wear): CRA e GDPR}
\label{sec:pulsek}
% =====================================================================

\begin{nota}
Le slide della simulazione non sono materia d'esame, ma questo tipo di esercizio (notifica alle
autorità) è molto probabile all'esame. Il caso aggiunge un livello di complessità in più: i dati
riguardano \textbf{minori}. Lo strumento \emph{CRA--GDPR Incident Desk} usato in aula è disponibile
sulla piattaforma del corso per esercitarsi.
\end{nota}

\subsection{Scenario}

\begin{itemize}
  \item \textbf{Volta Wear S.r.l.}, Torino (stabilimento principale in Italia), fondata nel 2019,
    48 dipendenti, fatturato 2025 di \texteuro{}9,4 milioni: \textbf{piccola impresa}.
  \item \textbf{Prodotto}: Pulse~K, smartwatch GPS per bambini di 5--12 anni, sul mercato UE da
    settembre 2024. Vendite tramite web shop proprio e retailer, tra cui ElektroMax GmbH (Germania).
  \item \textbf{182.000 orologi attivi}: Italia 118k, Francia 36k, Germania 17k, Spagna 11k.
\end{itemize}

\paragraph{Tre componenti, un solo prodotto.}
\begin{itemize}
  \item \textbf{Pulse K watch}: GPS, tasto SOS, messaggi vocali con contatti approvati;
  \item \textbf{Pulse Cloud}: progettato e gestito da Volta Wear su server UE in affitto; collega
    orologio e app;
  \item \textbf{Pulse Family app}: i genitori vedono la posizione live, impostano le zone ``Casa''
    e ``Scuola'', inviano messaggi vocali.
\end{itemize}

\paragraph{Dati trattati.}
\begin{center}
\begin{tabularx}{0.9\linewidth}{lXl}
\toprule
\textbf{Soggetto} & \textbf{Dato} & \textbf{Conservazione} \\ \midrule
Bambino  & Nome, anno di nascita, foto (facoltativa) & Durata dell'account \\
Bambino  & Posizione live e storico posizioni        & 30 giorni \\
Bambino  & Zone sicure ``Casa'' e ``Scuola''         & Durata dell'account \\
Bambino  & Messaggi vocali                           & 7 giorni \\
Genitore & Nome, email, cellulare                    & Durata dell'account \\ \bottomrule
\end{tabularx}
\end{center}
Non presenti nel cloud: dati di pagamento (gestiti dal payment provider) e password (salvate solo
come hash).

\subsection{La postura di sicurezza prima dell'incidente}

\begin{itemize}
  \item \textbf{Team}: un CTO e 9 sviluppatori, nessun addetto dedicato alla sicurezza.
  \item \textbf{Segnalazione vulnerabilità}: indirizzo \texttt{security@voltawear.eu} e una policy di
    due righe sul sito.
  \item \textbf{Componenti}: nessun elenco del software di terze parti incluso nel prodotto (nessuna SBOM).
  \item \textbf{Notti e weekend}: un solo sviluppatore reperibile; nessun legale né DPO reperibile;
    nessuna procedura di emergenza.
  \item \textbf{Customer care}: esternalizzato a CallPro S.r.l., con script fissi.
  \item \textbf{Privacy}: DPO esterno; \textbf{DPIA svolta al lancio nel 2024 e mai aggiornata}.
  \item \textbf{Formazione}: un e-learning generico di 45 minuti all'anno, uguale per tutti.
  \item \textbf{Ticket di aprile 2026}: ``aggiornare la libreria di login TokenGate, fix di sicurezza
    disponibile da marzo'', rinviato a dopo la campagna \emph{back to school}.
\end{itemize}

\begin{esame}
\textbf{DPIA non aggiornata.} Alla domanda ``esiste una DPIA?'' la risposta formalmente è sì, ma una
DPIA mai aggiornata può valere quasi quanto non averla. Averla resta meglio che non averla: tutto
dipende da quanto è cambiato il software dal 2024. Dettagli di questo tipo compaiono sia nei casi
reali sia nelle tracce d'esame.
\end{esame}

\subsection{Cronologia dell'incidente (i sei \emph{inject})}

La falla nel cloud era sfruttata da \textbf{mercoledì 4 novembre 2026, ore 21:00}.

\begin{center}
\small
\begin{tabularx}{\linewidth}{c p{2.6cm} X}
\toprule
\textbf{\#} & \textbf{Quando} & \textbf{Cosa succede} \\ \midrule
1 & Gio 5/11, 23:10 & Le richieste di posizione salgono a 40 volte il volume abituale, da centinaia
  di IP diversi. Lo sviluppatore reperibile blocca alcuni range IP e chiude il ticket come
  ``bot scraping, nessun impatto''. Nessun altro viene informato. \\ \midrule
2 & Ven 6/11, 08:30--11:05 & CallPro riceve 23 chiamate in due ore e mezza: gli orologi ricevono
  messaggi vocali da adulti sconosciuti. Gli operatori seguono lo script (riavviare l'orologio,
  cambiare il PIN). Alle 11:05 un team leader scrive al customer-care manager di Volta Wear;
  nessuno avvisa CTO o DPO. \\ \midrule
3 & Ven 6/11, 12:20 & Un ricercatore scrive a \texttt{security@}: cambiando il numero dell'orologio
  nella richiesta dell'app, qualsiasi utente autenticato vede la posizione live di qualsiasi bambino
  e può inviargli messaggi vocali. La falla è già sfruttata: un canale Telegram (``KidRadar'') vende
  la posizione live di qualsiasi Pulse~K a \texteuro{}15 a ricerca. Il ricercatore allega screenshot
  e annuncia la pubblicazione tra 30 giorni. \\ \midrule
4 & Ven 6/11, 16:00 & Riunione di crisi. I log confermano \textbf{41.250 orologi interrogati},
  \textbf{9.870 con posizione esposta} (posizione live, storico di 30 giorni, zone) e
  \textbf{37 bambini contattati} da sconosciuti. Cause: libreria TokenGate obsoleta (login
  falsificabile) e mancato controllo di \emph{ownership} (il cloud non verificava che l'orologio
  richiesto appartenesse all'account). CVSS 9.1 (critico). Il CTO propone di disattivare la
  condivisione della posizione; il CEO rifiuta: ``i genitori ci contano, non allarmiamo nessuno prima
  di lunedì, sistemiamo in silenzio''. \\ \midrule
5 & Sab 7/11 & Il retailer tedesco chiede se deve sospendere le vendite; una giornalista di
  \emph{Le Parisien} chiede un commento entro le 18:00; il marketing pubblica su Instagram:
  ``Pulse K è sicuro. Nessun dato è stato compromesso''. \\ \midrule
6 & Dom 8/11, 20:00 & Fix sul cloud: controllo di ownership aggiunto, TokenGate aggiornato, tutte le
  sessioni chiuse. Lunedì 9/11 alle 09:00 esce l'app 4.2. KidRadar resta online e ora vende gli
  \emph{storici} delle posizioni. \\ \bottomrule
\end{tabularx}
\end{center}

\subsection{Lezioni emerse dai singoli inject}

\paragraph{Inject 1 -- Sicurezza contro privacy.} Il riflesso tipico del tecnico di cybersicurezza è
``c'è un attacco, lo blocco''. Manca la sensibilità privacy: ogni attacco può avere impatto sui dati
personali, soprattutto quando riguarda minori.

\paragraph{Inject 2 -- Informazioni che non risalgono.} Il call center aveva segnali chiari, ma
nessuno ha avvisato DPO o CTO. Domanda aperta del debrief: ciò che sapeva un operatore del call
center esternalizzato vale come conoscenza dell'azienda?

\paragraph{Inject 3 -- Ricercatori e \emph{bug bounty}.} In alcuni casi l'attività è legale ed etica.
Diventa problematica quando la segnalazione si accompagna a richieste di denaro, che di fatto
equivalgono a un riscatto.

\paragraph{Inject 4 -- Bilanciamento di interessi.} Il CEO punta su reputazione e continuità del
servizio, ma non coglie che le conseguenze di un incidente gestito male possono essere molto più
gravi. ``Sistemare in silenzio'' è incompatibile con un orologio normativo da 24 ore.

\paragraph{Inject 5 -- Rischio reputazionale e comunicazione.} Esperienza maturata dal 2018 a oggi:
\begin{itemize}
  \item i clienti ricevono ormai spesso comunicazioni di data breach (es.\ il caso Trenitalia) e,
    soprattutto i più giovani, le accettano come normali;
  \item nel 2026 il rischio reputazionale di \emph{comunicare} un breach esiste ma raramente fa perdere
    clienti in modo significativo; il rischio vero è \textbf{non comunicare correttamente,
    minimizzare o mentire};
  \item i dati finiranno comunque sul dark web: la verità emergerà;
  \item ``chiedere scusa nel modo giusto'' è meglio che occultare;
  \item se i clienti sono già stati informati, con la stampa si è in posizione di forza: si può dire
    di aver già avvisato gli interessati e rifiutare dettagli coperti da indagini in corso;
  \item il post Instagram (``nessun dato compromesso'') contraddice i fatti noti all'azienda.
\end{itemize}
È compito del consulente spiegare all'alta direzione perché la trasparenza conviene.

\paragraph{Inject 6 -- Il fix fa partire un orologio CRA.} Si tratta del termine di 14 giorni per il
\emph{final report} sulla vulnerabilità (art.~14(2)(c)).

\subsection{Task 1 -- Qualificazione del caso ai sensi del CRA}

\begin{itemize}
  \item \textbf{È un prodotto con elementi digitali?} Sì. Pulse Cloud rientra come soluzione di
    elaborazione dati da remoto (\emph{remote data processing}) del prodotto.
  \item \textbf{Si applica un regime di esclusione?} No: non è un dispositivo medico, un veicolo,
    un prodotto dell'aviazione civile o dell'equipaggiamento marittimo, non è sviluppato
    esclusivamente per sicurezza nazionale o difesa, non è un pezzo di ricambio identico.
  \item \textbf{Ruolo}: Volta Wear è il \textbf{fabbricante}.
  \item \textbf{Categoria}: i wearable per bambini sono \textbf{prodotti importanti di Classe~I}
    (Allegato III).
  \item \textbf{Date}: l'evento (novembre 2026) cade dopo l'11/9/2026 (segnalazione già applicabile)
    e prima dell'11/12/2027 (applicazione integrale). Quindi \textbf{l'art.~14 si applica,
    l'Allegato~I non ancora}.
  \item \textbf{Dimensione}: piccola impresa; per le micro e piccole imprese l'art.~64(10) prevede una
    regola sanzionatoria specifica (vedi il riquadro \emph{Da verificare} in fondo alla lezione).
  \item \textbf{Verdetto}: c'è evidenza affidabile che un attore malevolo ha sfruttato la
    vulnerabilità, e l'incidente compromette la capacità del prodotto di proteggere dati sensibili o
    importanti. Sono soddisfatti entrambi i presupposti: \textbf{vulnerabilità attivamente sfruttata}
    e \textbf{incidente grave}, cioè \textbf{due binari di notifica CRA}.
  \item C'è inoltre, senza dubbio, un \textbf{data breach ai sensi del GDPR}.
\end{itemize}

\subsection{Task 2 -- Gli orologi: CRA e GDPR}

\paragraph{Consapevolezza ai fini del CRA.}
\begin{itemize}
  \item \textbf{Vulnerabilità attivamente sfruttata}: in aula (professore e Gianluca concordi) si è
    adottata la lettura prudente, che fa coincidere la consapevolezza con la prima evidenza
    disponibile all'azienda: l'alert di giovedì 5/11 alle 23:10. Lettura alternativa: venerdì 6/11
    alle 12:20, quando l'email del ricercatore fornisce l'evidenza affidabile dello sfruttamento.
  \item \textbf{Incidente grave}: venerdì 6/11 alle 16:00 (riunione di crisi, i log confermano).
  \item \textbf{Mitigazione disponibile}: domenica 8/11 alle 20:00.
\end{itemize}

\paragraph{Consapevolezza ai fini del GDPR.}
Il titolare è consapevole quando ha un ragionevole grado di certezza che \textbf{dati personali}
siano stati compromessi. Qui le opzioni (5 o 6 novembre) distano circa un giorno e la scelta è poco
delicata. Diventa invece delicatissima quando, per esempio, l'attacco avviene il 1° e il CEO viene
informato l'8:
\begin{itemize}
  \item per sostenere una data di consapevolezza posticipata servono \textbf{evidenze documentali} del
    momento in cui il titolare è stato informato: la \textbf{tracciabilità} nella gestione
    dell'incidente è essenziale;
  \item senza evidenze si rischia una \textbf{dichiarazione falsa} al Garante, punita con la
    reclusione da sei mesi a tre anni;
  \item \textbf{posizione di Gianluca}: l'orologio GDPR parte quando si è consapevoli della
    compromissione di dati personali. L'attacco può essere in corso da giorni, settimane o mesi
    (ricognizione, movimento laterale, ecc.): queste fasi fanno parte dell'incidente per CRA, DORA e
    NIS2, ma non rilevano per l'orologio GDPR finché non toccano dati personali;
  \item \textbf{aggiunta del professore}: se l'attacco era evidente e qualcuno (es.\ il CISO) l'ha
    nascosto pensando di gestirlo da solo, è più prudente non far partire l'orologio dalla data in
    cui è stato informato il CEO, perché un ritardo così lungo non sarebbe credibile. Un conto è un
    attacco occulto rimasto ignoto per mesi, altro è un attacco evidente tenuto nascosto.
\end{itemize}

\begin{center}
\textbf{Tabella delle scadenze}\par\smallskip
\small
\begin{tabularx}{\linewidth}{l X l l}
\toprule
\textbf{Binario} & \textbf{Obbligo} & \textbf{Decorrenza} & \textbf{Scadenza} \\ \midrule
CRA  & Early warning vulnerabilità (14(2)(a))     & Gio 5/11 23:10       & Ven 6/11 23:10 \\
CRA  & Notifica vulnerabilità (14(2)(b))          & Gio 5/11 23:10       & \textbf{Dom} 8/11 23:10 \\
CRA  & Final report vulnerabilità (14(2)(c))      & Dom 8/11 20:00 (fix) & \textbf{Dom} 22/11 20:00 \\
CRA  & Early warning incidente grave (14(4)(a))   & Ven 6/11 16:00       & \textbf{Sab} 7/11 16:00 \\
CRA  & Notifica incidente (14(4)(b))              & Ven 6/11 16:00       & Lun 9/11 16:00 \\
CRA  & Final report incidente (14(4)(c))          & Notifica 14(4)(b)    & $\approx$ Mer 9/12 \\
CRA  & Informare gli utenti (14(8))               & Consapevolezza       & Tempestivamente \\
GDPR & Notifica al Garante (art.~33), lettura prudente & Gio 5/11 23:10  & \textbf{Dom} 8/11 23:10 \\
GDPR & Notifica al Garante (art.~33), lettura tardiva  & Ven 6/11 16:00  & Lun 9/11 16:00 \\
GDPR & Comunicazione alle famiglie (art.~34)      & Consapevolezza       & Senza ingiustificato ritardo \\
\bottomrule
\end{tabularx}
\end{center}

\noindent Canale CRA: Single Reporting Platform ENISA, verso il CSIRT dello Stato dello stabilimento
principale (Italia). L'early warning deve indicare gli Stati membri in cui il prodotto è messo a
disposizione (Italia, Francia, Germania, Spagna).

\noindent\textbf{Quanto costa ``aspettare lunedì''}: entro lunedì mattina sono già scaduti l'early
warning sulla vulnerabilità, la notifica della vulnerabilità, l'early warning sull'incidente grave e,
nella lettura prudente, anche le 72 ore GDPR.

\subsection{Task 3 -- Severità del breach (metodologia ENISA 2013)}

\[
SE = DPC \times EI + CB
\]
Scale e valori sono riepilogati nella scheda di ripasso (Capitolo~\ref{cap:schede}).

\paragraph{Gruppo A -- 31.380 famiglie (solo dati di profilo).}
Nome, anno di nascita, foto del bambino; nome, email, cellulare del genitore.
DPC = 1 (dati semplici), EI = 1 (massima), CB = 0,5 (riservatezza, destinatari ignoti) + 0,5
(intento malevolo: vendita su Telegram).
\[
SE_A = 1 \times 1 + 1 = 2 \quad\Rightarrow\quad \text{media}
\]

\paragraph{Gruppo B -- 9.870 famiglie (profilo e posizione).}
Profilo, posizione live, storico di 30 giorni, zone ``Casa'' e ``Scuola''; 37 bambini contattati.
DPC = 2 (dati comportamentali), EI = 1, CB = 0,5 + 0,5.
\[
SE_B = 2 \times 1 + 1 = 3 \quad\Rightarrow\quad \text{alta (parte bassa della fascia)}
\]

\paragraph{Critica alla metodologia.} Il professore la definisce ``folle'' su questo punto: la
metodologia ENISA, con un'impostazione pensata per il settore finanziario, \textbf{non distingue tra
minori e adulti}. Gianluca precisa che, in senso stretto, i dati di profilo sono DPC~1, ma il valore
può essere \textbf{aumentato motivatamente} in base al numero di interessati e alla loro
vulnerabilità (\emph{accountability}).

\paragraph{Decisioni.}
\begin{itemize}
  \item \textbf{Art.~33}: notifica al Garante \textbf{obbligatoria}. Il Garante italiano è l'autorità
    capofila (stabilimento principale a Torino, trattamento transfrontaliero: art.~56).
  \item \textbf{Art.~34}: comunicazione agli interessati del Gruppo~B \textbf{sì}, anche se il punteggio
    (3) consentirebbe di discuterne e il CEO probabilmente spingerebbe per evitarla. Le ragioni:
    \begin{itemize}
      \item trasparenza;
      \item interessati estremamente vulnerabili: i dati possono interessare criminali, inclusi pedofili;
      \item secondo Gianluca, lavoratori, malati e bambini sono categorie particolarmente seguite
        dalle autorità di controllo, quindi il rischio va valutato al livello più alto possibile.
    \end{itemize}
  \item \textbf{A chi si comunica}: ai \textbf{genitori} (titolari della responsabilità genitoriale),
    non ai bambini.
\end{itemize}

\subsection{Task 4 -- Piano di remediation (da svolgere a casa)}

Per mancanza di tempo il task è rimasto da svolgere individualmente con l'Incident Desk.
Il piano va articolato su quattro assi:
\begin{itemize}
  \item \textbf{tecnico}: chiudere le due falle e fare in modo che il prossimo componente obsoleto
    venga individuato in giorni, non in mesi;
  \item \textbf{organizzativo}: chi decide alle 23:10 di un giovedì, e quali template devono già esistere;
  \item \textbf{formazione}: almeno tre azioni, ciascuna con destinatari, contenuto e KPI, collegate a
    un inject;
  \item \textbf{comunicazione}: famiglie, tutti gli utenti, retailer, stampa, ricercatore, rettifica
    del post Instagram.
\end{itemize}

\begin{esame}
Serve un \textbf{piano di risposta agli incidenti integrato} che coordini contemporaneamente CRA,
GDPR e, dove applicabile, NIS2 (tema della prossima lezione).
\end{esame}

% =====================================================================
\section{Prossimo laboratorio (venerdì): il ruolo dell'auditor}
% =====================================================================

Presentato da Gianluca. Cambia il ruolo: non si gestisce più l'incidente dall'interno, si è
\textbf{auditor esterno}.

\paragraph{Scenario.}
\begin{itemize}
  \item \textbf{Sestante Mobility}, \textbf{media impresa} (non micro né piccola: si applicano tutti
    gli obblighi previsti per le medie imprese).
  \item Prodotto: box telematico installato su veicoli di flotte aziendali, collegato a una porta del
    veicolo da cui estrae dati, con piattaforma cloud (SaaS) e app mobile.
  \item Clienti: \emph{fleet operator} in diversi Paesi.
  \item Dati raccolti: viaggi, eventi di guida e un punteggio del conducente calcolato dal dispositivo.
    Il conducente ha nome e cognome, quindi si tratta di \textbf{dati personali}.
\end{itemize}

\paragraph{Eventi (date nel futuro).}
\begin{itemize}
  \item \textbf{Agosto 2027}: un ricercatore segnala che il firmware accetta manomissioni.
  \item \textbf{Settembre 2027}: nei dispositivi ricondizionati (ritirati da ex clienti e riassegnati
    a nuovi clienti) restano dati del cliente precedente.
  \item \textbf{Novembre 2027}: un conducente scopre di essere tracciato 24 ore su 24, anche fuori
    dall'orario di lavoro (i veicoli restano in uso al conducente), e il cliente non riesce a
    disattivare il tracciamento.
\end{itemize}

\paragraph{Domande a cui rispondere.}
\begin{itemize}
  \item Si può mantenere il prodotto sul mercato, in particolare le unità vendute dopo il
    \textbf{12 dicembre 2027}?
  \item Come gestire un cliente che rinnova parte dei dispositivi entro il \textbf{24 aprile 2028}?
    Le unità sono conformi?
  \item Un nuovo dispositivo non ancora sul mercato: è valutabile? Senza evidenze, l'auditor deve
    dichiararlo fuori perimetro o non valutabile.
  \item Che ruolo ha Sestante ai fini GDPR: titolare, responsabile o entrambi?
\end{itemize}

\begin{esame}
\textbf{Regole metodologiche dell'audit:}
\begin{itemize}
  \item le date di cut-off si applicano alla \textbf{singola unità} immessa sul mercato, non al
    modello: unità dello stesso modello possono essere conformi o no a seconda della data di immissione;
  \item \textbf{una fonte} = asserzione, da prendere con riserva; \textbf{due fonti concordanti} =
    fatto accertato (evidenza);
  \item \textbf{un certificato ISO non implica conformità CRA né GDPR}: il GDPR è una norma giuridica,
    la ISO uno standard tecnico di sicurezza delle informazioni;
  \item la scala di severità dell'audit è diversa da quella ISO (spiegata nel playbook);
  \item il perimetro va preso sul serio: \textbf{non inventare né presumere nulla}. Ciò che non è
    nelle evidenze non esiste.
\end{itemize}
\end{esame}

\paragraph{Materiale e consegna.} 16 documenti brevi (poche frasi ciascuno) e le trascrizioni di due
brevi interviste (tra cui il CTO), con un ordine di analisi suggerito: prima disclosure, configurazione
e impostazioni di default del prodotto, poi le informazioni di secondo livello. Briefing e playbook
vengono caricati in anticipo; la piattaforma online è attiva solo il giorno del laboratorio.
Deliverable: un testo breve, da compilare online e caricare come file di testo secondo le istruzioni.

% ---------------------------------------------------------------------
\section{Domande aperte e punti da verificare}

\begin{verifica}
\begin{itemize}
  \item \textbf{Data di consapevolezza CRA}: la trascrizione fa dire alla classe ``4 novembre'', ma il
    4 è il giorno di inizio dello sfruttamento, di cui l'azienda non sapeva nulla. Qui si è usato
    l'alert del 5/11 alle 23:10. Confermare con i propri appunti.
  \item \textbf{Tabella delle scadenze}: in aula il calcolo non è stato completato; le date sono
    calcolate applicando art.~14 CRA e art.~33 GDPR alle decorrenze scelte.
  \item \textbf{Art.~64(10) CRA}: per micro e piccole imprese non si applicano sanzioni per il solo
    mancato rispetto del termine di 24 ore dell'early warning. Verificare sul testo ufficiale.
  \item \textbf{Art.~14(8) CRA}: l'obbligo di informare gli utenti non ha un termine in ore.
    Verificare sul testo ufficiale.
  \item Cognome del co-docente Gianluca (non ricavabile con certezza dalla trascrizione).
\end{itemize}
\end{verifica}
